% !TEX root = ../aa_SoM_Chapters.tex

%% HARRIS: Chapter 9
%\xdef\Isectiontitle{Statistical Mechanics} 
%%
%\xdef\IaSubsectiontitle{A Simple Thermodynamic System} 
%\xdef\IbSubsectiontitle{Entropy and Temperature}
%\xdef\IcSubsectiontitle{Einstein's solid}
%\xdef\IdSubsectiontitle{Boltzmann \& Maxwell–Boltzmann distributions}
%\xdef\IeSubsectiontitle{Classical Averages}
%
%\xdef\IfSubsectiontitle{Quantum Distributions} 
%\xdef\IgSubsectiontitle{The Quantum Gas} 
%\xdef\IhSubsectiontitle{Photon Gas}
%\xdef\IiSubsectiontitle{The Laser}
%\xdef\IjSubsectiontitle{Specific Heats} 
%\xdef\IkSubsectiontitle{Debye Model} 

% ö = \%c3\%b6
% ä = \%C3\%A4

\xdef\sectiontitle{\IVsectiontitle} %aiheuttama

%%%%%%%%%%%%%%%
\begin{frame}[label=4_7_a]
\frametitle{\texorpdfstring{\culine[\sectiontitleulinecolor]{\sectiontitle}}{\sectiontitle} }

%\vspace{\chapterTOCvksip}%
\begin{itemize}[leftmargin=0.5cm,itemsep=\chapterTOCitemsep]
    \item[4.1] \hyperlink{4_1_a}{\IVaSubsectiontitle}
    \item[4.2] \hyperlink{4_2_a}{\IVbSubsectiontitle}	
    \item[4.3] \hyperlink{4_3_a}{\IVcSubsectiontitle}	
    \item[4.4] \hyperlink{4_4_a}{\IVdSubsectiontitle}
    \item[4.5] \hyperlink{4_5_a}{\IVeSubsectiontitle}
    \item[4.6] \hyperlink{4_6_a}{\IVfSubsectiontitle}
    \item[4.7] \hyperlink{4_7_a}{\CDAlert[blue]{\IVgSubsectiontitle}}
\end{itemize}

%\endofslide 
\end{frame}
%%%%%%%%%%%%%%%

\xdef\sectiontitle{\IVgSubsectiontitle} 
\xdef\sectiontitle{GUT and Cosmology} 


%%%%%%%%%%%%%%%
\begin{frame}
\frametitle{\texorpdfstring{\culine[\sectiontitleulinecolor]{\sectiontitle}}{\sectiontitle} }
\framesubtitle{Fundamental Particles \& background radiation}

\begin{itemize}
	\item In recent years\appear{, \CDAlert[red]{high-energy physics}} \appear{and \CDAlert[blue]{cosmology} have approached each other}
	
	\item This is because the theories describing the \CDAlert[fuchsia]{fundamental interactions} \appear{are linked to the early history of the universe}\appear{, right after the \CDAlert[fuchsia]{Big Bang}}\appear{, and how the fundamental particles and the material in the universe had its origin}

	\item For some time after the Big Bang \appear{the energies were quite high}\appear{, exceeding the atomic building energies}\appear{, the material consisted of atoms and plasma of electrons and positive ions}
	
	\item Gradually the temperature dropped\appear{, and presently the \Alert[tuniblue]{cosmic background radiation} is of the order of $2.7 \mathrm{\;K}$}\appear{, which value is in agreement with the present theories on the \CDAlert[blue]{universe expanding after Big Bang}}
\end{itemize}


\endofslide 
%\endofsection
\end{frame}
%%%%%%%%%%%%%%%

%%%%%%%%%%%%%%%
\begin{frame}
\frametitle{\texorpdfstring{\culine[\sectiontitleulinecolor]{\sectiontitle}}{\sectiontitle} }
\framesubtitle{Expanding universe}

\begin{itemize}
	\item \CDAlert[fuchsia]{Expanding universe} can be qualitatively described by a simple model. \appear{Let's consider a projectile}\appear{, or a satellite $(m)$} \appear{launched directly upwards from a planet ($M$).} \appear{Assuming that $r$ is the distance of the satellite from the planet's center}\appear{, and $\rho$ is the density of the planet}\appear{, the energy is given by:}
\[
E  \equal \frac{1}{2} \appear{m v^2}   \plus  \appear{E_{pot} } \equal \frac{m}{2}  \appear{\textrm{((}} \frac{d r}{d t} \appear{\textrm{))}^2}   \minus \fraca{G m M}{r}  \equal \frac{m}{2}  \appear{\textrm{((}} \frac{d r}{d t} \appear{\textrm{))}^2}   \minus \frac{G m}{r} \appear{\rho} \frac{4}{3} \appear{\pi r^{3}}
\]
\item[] If \CDAlert[red]{$E<0$}\appear{, then the projectile will return back to the planet }
\item[] If \CDAlert[blue]{$E>0$}\appear{, then the \CDAlert[blue]{projectile can escape}}

\item Rearranging, we get the equation:
\[
 \textrm{((} \frac{d r / d t}{r} \appear{\textrm{))}^2}   \equal \frac{8 \pi G \rho}{3}  \plus \frac{2 E / m}{r^{2}}
\]
\end{itemize}

\endofslide 
%\endofsection
\end{frame}
%%%%%%%%%%%%%%%

%%%%%%%%%%%%%%%
\begin{frame}
\frametitle{\texorpdfstring{\culine[\sectiontitleulinecolor]{\sectiontitle}}{\sectiontitle} }
\framesubtitle{Friedmann equation}

\begin{itemize}
	\item Quite \Alert[fuchsia]{analogous to the escaping projectile}\appear{, we can assume that \CDAlert[fuchsia]{universe} is an expanding sphere}\appear{, with size $R=R(t)$}\appear{, actually called a \CDAlert[red]{scale factor}} 
	
	\item Expansion is described by \CDAlert[blue]{Friedmann equation}:
\[
 \textrm{((} \frac{d R / d t}{R} \appear{\textrm{))}^2}  \equal \frac{8 \pi G \rho}{3}  \minus \frac{K}{R^{2}}
\]
\appear{The constant $K$ is analogous to the energy of the projectile}

\item If \Alert[red]{$K>0$:} \appear{positively curved universe}\appear{, expansion stops}\appear{, and universe will crunch back to $R=0$}

\item If \Alert[green]{$K=0$:} \appear{a flat universe}\appear{, expansion constant to infinite time $t \rightarrow \infty$ }
\item If \Alert[blue]{$K<0$:} \appear{negatively curved universe}\appear{, expansion will accelerate indefinitely}
\end{itemize}

\endofslide 
%\endofsection
\end{frame}
%%%%%%%%%%%%%%%

%%%%%%%%%%%%%%%
\begin{frame}
\frametitle{\texorpdfstring{\culine[\sectiontitleulinecolor]{\sectiontitle}}{\sectiontitle} }
\framesubtitle{Friedmann equation}

\begin{itemize}
	\item Next, we modify the density in the equation: $\appear{\rho} \rightarrow \appear{\rho_{M} /{R^{3}}}  \plus \appear{\rho_{\Lambda}}$ \appear{where $\rho_{M} /R^{3}$ is the present density of the universe} \appear{and $\rho_{\Lambda}$ is a factor connected to the \CDAlert[fuchsia]{cosmological constant}} \appear{(introduced later):}
\[
 \textrm{((} \frac{d R / d t}{R} \appear{\textrm{))}^2}   \equal \frac{8 \pi G \rho_{M}}{3 R^{3}}  \plus \frac{8 \pi G \rho_{\Lambda}}{3}  \minus \frac{K}{R^{2}}
\]
\appear{Finally, we rearrange the Friedmann equation to its final form:}
\[
 \textrm{((} \frac{d R / d t}{R} \appear{\textrm{))}^2}   \equal \frac{\Omega_{M}}{R^{3}}  \plus \appear{\Omega_{\Lambda}}  \minus \frac{K^{\prime}}{R^{2}} \quad \appear{\text { with: } R_{\text {present}}=1} \appear{\text { ~and~ }(d R / d t)_{\text {present}}=1}
\]
\item[] We have in essence \CDAlert[blue]{three tuning parameters} \appear{$\Omega_{M}$}\appear{, $\Omega_{\Lambda}$} \appear{and $K^{\prime}$}
\end{itemize}

\endofslide 
%\endofsection
\end{frame}
%%%%%%%%%%%%%%%

%%%%%%%%%%%%%%%
\begin{frame}
\frametitle{\texorpdfstring{\culine[\sectiontitleulinecolor]{\sectiontitle}}{\sectiontitle} }
\framesubtitle{Einstein's cosmological constant}

\begin{itemize}
	\item On the tuning parameters: %\vspace{1mm}
	\begin{itemize}
	\item $K^{\prime}=0$\appear{, assuming a flat universe,
having nearly constant expansion}
\item $\Omega_{M}=1$ \appear{presently $\Leftrightarrow$ critical density}
\item $\Omega_{\Lambda} \approx 0$

\end{itemize}
\item Using above assumptions we get
\end{itemize}

% 6.5 cm = midway, 10 cm = 3/4 
%\vspace{2.5mm}
\begin{minipage}{7.5cm}
\begin{itemize}
	\item[] $$
	\begin{aligned}
	\appear{\textrm{((}} \frac{d R / d t}{R} \appear{\textrm{))}^2} & \equal  \frac{1}{R^{3}} \Rightarrow \appear{\nint_{1}^{R}} \appear{R^{1 / 2}} \appear{\rmd R}  \equal \appear{\nint_{1}^{t} d t} \\
\appear{\Rightarrow \quad R(t)} & \equal  \appear{\textrm{[[} \Frac{3}{2}} \appear{( t-\frac{1}{3} )} \appear{\textrm{]]}^{2 / 3}}
	\end{aligned}
$$
\item[] $\Omega_{\Lambda}=$ \CDAlert[red]{cosmological constant}\appear{, potentially determining whether the expansion of the universe will be \Alert[red]{constant or not}} \appear{(see Fig 17)}
\end{itemize}
\end{minipage}
%\vspace{2.5mm}

\appear{\xdef\loc{south east}
\begin{tikzpicture}[remember picture,overlay] % anchor=south
    \node (A) [yshift=-0.15*\figYshift,xshift=\figXshift, anchor=\loc] at (current page.\loc) {\includegraphics[page=1,width=6.5cm]{Figures_booklet/SOM_11_Fundamental_particles_figures/SOM_Fundamental_particles_Figure_17.png}}; %scale=\figscale. % 8cm max hei
\end{tikzpicture}}

\endofslide 
%\endofsection
\end{frame}
%%%%%%%%%%%%%%%

%%%%%%%%%%%%%%%
\begin{frame}
\frametitle{\texorpdfstring{\culine[\sectiontitleulinecolor]{\sectiontitle}}{\sectiontitle} }
\framesubtitle{Chronology of the Universe}

\begin{itemize}
	\item Earlier, we looked into proton–proton cycle occurring within stars

\item In fact, history of the whole universe has some similarities

\item At the present knowledge \appear{((called the standard model of the history of the universe))}\appear{, the temperature of the universe at a time of $t=10^{-13} \mathrm{\;s}$ (after Big Bang) was about $10^{32} \mathrm{\;K}$}\appear{, and the average energy per particle was:} \vspace{-4mm}
\[
\hspace{2cm} \appear{E} \approx \LP \appear{10^{-13}} \frac{\mathrm{\;GeV}}{\mathrm{K}} \,\RP  \appear{10^{32} \mathrm{\;K}}  \equal \appear{10^{19} \mathrm{\;GeV}}
\]

\item This time has been suggested to be the time for the transition from the \CDAlert[red]{Theory of Everything (TOE)} \appear{to \CDAlert[blue]{Grand Unified Theory (GUT)}}

\item At $t=10^{-35} \mathrm{\;s}$ the temperature had decreased to about $10^{27} \mathrm{\;K}$\appear{, the energy being in the order of $10^{14} \mathrm{\;GeV}$.} \appear{\Alert[fuchsia]{At this energy the strong force separated from electroweak force}}
\end{itemize}

\endofslide 
%\endofsection
\end{frame}
%%%%%%%%%%%%%%%

%%%%%%%%%%%%%%%
\begin{frame}
\frametitle{\texorpdfstring{\culine[\sectiontitleulinecolor]{\sectiontitle}}{\sectiontitle} }
\framesubtitle{Chronology of the Universe}

\semismall
\begin{itemize}[itemsep=1.7mm]
	\item At $\mathrm{t}=10^{-32} \mathrm{\;s}$ \appear{the temperature had decreased to ${\sim}10^{25} \mathrm{\;K}$} \appear{and the universe was a mixture of \CDAlert[fuchsia]{quarks}}\appear{, \CDAlert[blue]{leptons}} \appear{and the \CDAlert[red]{mediating bosons}}
	
	\item At $\mathrm{t}=10^{-6} \mathrm{\;s}$ \appear{the temperature had decreased to about $10^{13} \mathrm{\;K}$} \appear{and the typical energy was about $1 \mathrm{\;GeV}$} \appear{(comparable to rest mass of a nucleon).} \appear{\CDAlert[fuchsia]{Quarks began binding together to from nucleons}}
	
	\item At about $\mathrm{t}=1 \mathrm{\;s}$\appear{, the ratio of protons to neutrons was determined by the Boltzmann distribution factor $e^{-\Delta E / k_{B} T}$} \appear{where $\Delta E$ is the difference between the neutron and proton rest energies} $\appear{\Delta E}  \equal \appear{1.29 \mathrm{\;MeV}}$
	
	\item At $T \Approx10^{10} \mathrm{\;K}$ \appear{the distribution gives about $4.5$ times as many protons as neutrons}
	
	\item The formation of deuterons started at \CDAlert[red]{$\mathrm{t}\Approx225 \mathrm{\;s}$}\appear{, marking the beginning of the period of nuclei formation}\appear{, i.e. \CDAlert[red]{Big Bang nucleosynthesis (BBN)}} \appear{or primordial nucleosynthesis}
	
	\item Further nucleosynthesis \appear{(\CDAlert[blue]{stellar nucleosynthesis})} \appear{occurred much later at \CDAlert[blue]{$t=10^{13} \mathrm{\;s}$}} \appear{when $T \sim 3000 \mathrm{\;K}$}\appear{, and average energy was few tenths of an electron volt}
	
\end{itemize}

\endofslide 
%\endofsection
\end{frame}
%%%%%%%%%%%%%%%

%%%%%%%%%%%%%%%
\begin{frame}
\frametitle{\texorpdfstring{\culine[\sectiontitleulinecolor]{\sectiontitle}}{\sectiontitle} }
\framesubtitle{Timeline 1}

\xdef\loc{south}
\begin{tikzpicture}[remember picture,overlay] % anchor=south
    \node (A) [yshift=-0.25*\figYshift,xshift=\figXshift, anchor=\loc] at (current page.\loc) {\includegraphics[page=1,height=8.5cm]{Figures_booklet/SOM_11_Fundamental_particles_figures/SOM_Fundamental_particles_Figure_18.png}}; %scale=\figscale. % 8cm max hei
\end{tikzpicture}

\endofslide 
%\endofsection
\end{frame}
%%%%%%%%%%%%%%%

%%%%%%%%%%%%%%%
\begin{frame}
\frametitle{\texorpdfstring{\culine[\sectiontitleulinecolor]{\sectiontitle}}{\sectiontitle} }
\framesubtitle{Timeline 2}

%\xdef\loc{south}
%\begin{tikzpicture}[remember picture,overlay] % anchor=south
%    \node (A) [yshift=-0.25*\figYshift,xshift=\figXshift, anchor=\loc] at (current page.\loc) {\includegraphics[page=1,width=9.5cm]{Figures_booklet/SOM_11_Fundamental_particles_figures/SOM_Fundamental_particles_Figure_1.png}}; %scale=\figscale. % 8cm max hei
%\end{tikzpicture}

\endofslide 
%\endofsection
\end{frame}
%%%%%%%%%%%%%%%

%%%%%%%%%%%%%%%
\begin{frame}
\frametitle{\texorpdfstring{\culine[\sectiontitleulinecolor]{\sectiontitle}}{\sectiontitle} }
\framesubtitle{Timeline 3}


%\xdef\loc{south}
%\begin{tikzpicture}[remember picture,overlay] % anchor=south
%    \node (A) [yshift=-0.25*\figYshift,xshift=\figXshift, anchor=\loc] at (current page.\loc) {\includegraphics[page=1,width=9.5cm]{Figures_booklet/SOM_11_Fundamental_particles_figures/SOM_Fundamental_particles_Figure_1.png}}; %scale=\figscale. % 8cm max hei
%\end{tikzpicture}

%\endofslide 
\endofsection
\end{frame}
%%%%%%%%%%%%%%%

